\clearpage
\section{Sustentabilidad}

El proyecto tiene como objetivo procesar los desechos vegetales y frutales generados en una casa, por lo cual se aprovechan para crear abono casero para plantas o que puede ser aprovechado para áreas verdes públicas. Esto en miras de los objetivos ambientales que se tienen dentro de la Ciudad de México.

También se abre como una opción que busca estar al alcance de la gente para que se integre en las tareas de reutilización de recursos, generando un dispositivo que sea accesible y que se pueda adaptar en una casa común dentro del Valle de México.

El diseño del prototipo está previsto para ser manipulado por jóvenes y adultos, debido a que el sistema de trituración se encuentra expuesto y el sistema de regulación de temperatura puede representar un riesgo para los niños. Sin embargo, se contemplaron las condiciones necesarias para que una persona externa pueda operar el dispositivo de manera segura y, al mismo tiempo, participar en el proceso de compostaje.

Con su implementación, se busca tener un impacto ecológico positivo al realizar el procesamiento y aprovechamiento de residuos orgánicos, lo cual puede generar un beneficio económico al ahorrar dinero por la compra de fertilizantes, así como contribuir a la eliminación de malos olores y alejar los desechos de las calles.

Este proyecto tiene como objetivo a largo plazo acercar los procesos de compostaje a la población civil, facilitando el aprovechamiento de los residuos orgánicos y contribuyendo a la reducción de su impacto ambiental.